\subsection{Интересные факты из жизни}
Количество остовных деревьев в полном графе $n^{n - 2}$. \\
Перебор подмасок: \text{s = (s - 1) \& m} \\
Код Грея: $n = n \land (n >> 1)$ \\

\lstinputlisting[language=Octave]{source/Mathematics/binary_input.cpp}

Пусть $p$ -- простое число. Тогда $(p - 1)! \equiv -1 \mod p$. \\
Если сравнение $a_nx^n+ ... + a_1x+a_0 \equiv 0 \mod p$ степени $n$ по простому модулю $p$ имеет более $n$ различных решений, то все коэффициенты $a_0, a_1, . . . , a_n$ кратны $p$. \\

\subsubsection{Радиус вписанной и описанной окружностей}

The radius of circumscribed circle of triangle
$R = \frac{a\cdot b \cdot c}{4\cdot S}$

The radius of circumscribed circle of regular polygon
$R = \frac{a}{2\cdot \sin (\frac{180^o}{n})}$

The radius of inscribed circle of triangle
$R = \frac{S}{p},\ p\ \text{--- half-perimeter.}$

\subsubsection {Объём и площадь геометрических фигур}
\begin{minipage}{0.48\textwidth}
	\begin{multicols}{2}
	\textbf {Sphere}:

	$$S = 4 \cdot \pi R^2;\ S = 2 \cdot \pi \cdot R\cdot h,\ V = \frac{4}{3} \pi \cdot R^3$$

	The surface area of the sphere sector: 
	$$S = 2 \cdot\pi \cdot R \cdot (R + \frac{r}{2} - \sqrt{R^2 - r^2}) $$
	
	The $V$ of sphere segment: 
	$$V = \frac{1}{3} \cdot \pi \cdot h^2 \cdot (3\cdot R - h)$$

	The $V$ of sphere sector: 
	$$V = \frac{2}{3} \cdot \pi \cdot R^2 \cdot h$$

	The $V$ of sphere layer: 
	$$V = \frac{1}{3} \cdot \pi \cdot h \cdot (R^2 + r^2 + \frac{1}{3} \cdot h^2)$$


	\textbf {Equilateral triangle}:

	$$S = \frac{\sqrt{3}}{4} \cdot a^2$$

	\textbf {Cone}:

	$$S = \pi \cdot r \cdot l + \pi \cdot r^2,\ V = \frac{h}{3} \cdot \pi \cdot r^2$$

	\textbf {Regular pyramid}:

	$$S = \frac{1}{2}\cdot P \cdot L,\ \text {$P$ --- perimeter, $L$ --- pyramid former.}$$

	$$V = \frac{1}{3}\pi \cdot R^2 \cdot H,\ V = \frac{1}{3}\pi \cdot h (R^2 + R\cdot r + r^2)$$

	\textbf{Regular tetrahedron}:

	$$V = \frac{a^3 \cdot \sqrt{2}}{12} $$

	\textbf{Pyramid}

	$$V = \frac{1}{3} \cdot S \cdot h$$

	V of truncated pyramid: $$V = \frac{1}{3} \cdot h \cdot(S_{down} + \sqrt{S_{down} \cdot S_{up}} + S_{up}) $$

	\textbf {Torus}

	$$V = 2 \cdot \pi^2 R \cdot r^2,\ S = 4 \cdot \pi^2 \cdot R \cdot r$$

	\end{multicols}
\end{minipage}


\subsubsection{Тригонометрия}
\begin{minipage}{0.48\textwidth}
	\begin{multicols}{2}
		$$\sin{\alpha} + \sin{\beta} = 2 \cdot \sin{\frac{\alpha + \beta}{2}} \cdot \cos{\frac{\alpha - \beta}{2}}$$
		$$\sin{\alpha} - \sin{\beta} = 2 \cdot \sin{\frac{\alpha - \beta}{2}} \cdot \cos{\frac{\alpha + \beta}{2}}$$

		$$\cos{\alpha} + \cos{\beta} = 2 \cdot \cos{\frac{\alpha + \beta}{2}} \cdot \cos{\frac{\alpha - \beta}{2}}$$
		$$\cos{\alpha} - \cos{\beta} = - 2 \cdot \sin{\frac{\alpha + \beta}{2}} \cdot \cos{\frac{\alpha - \beta}{2}}$$

		$$\sin{(\alpha + \beta)} = \sin{\alpha} \cdot \cos{\beta} + \cos{\alpha} \cdot \sin{\beta}$$
		$$\sin{(\alpha - \beta)} = \sin{\alpha} \cdot \cos{\beta} - \cos{\alpha} \cdot \sin{\beta}$$

		$$\cos{(\alpha + \beta)} = \cos{\alpha} \cdot \cos{\beta} - \sin{\alpha} \cdot \sin{\beta}$$
		$$\cos{(\alpha - \beta)} = \cos{\alpha} \cdot \cos{\beta} + \sin{\alpha} \cdot \sin{\beta}$$
	
		$$\sin{\alpha} \cdot \sin{\beta} = \frac{\cos{(\alpha - \beta)} - \cos{(\alpha + \beta)}}{2}$$

		$$\sin{\alpha} \cdot \cos{\beta} = \frac{\sin{(\alpha - \beta)} + \sin{(\alpha + \beta)}}{2}$$

		$$\cos{\alpha} \cdot \cos{\beta} = \frac{\cos{(\alpha - \beta)} + \cos{(\alpha + \beta)}}{2}$$
	\end{multicols}
\end{minipage}
\newline
\newline

For even $n$:

$$cos^n \alpha=\frac {C_\frac n 2^n}{2^n}+\frac1{2^{n-1}} \cdot \sum_{k=0}^{\frac{n}{2} - 1} C_k^n \cdot \cos{((n-2k) \cdot \alpha)}$$

$$sin^n \alpha=\frac {C_\frac n 2^n}{2^n}+\frac1{2^{n-1}} \cdot \sum_{k=0}^{\frac{n}{2} -1} (-1)^{\frac n 2 -k} \cdot C_k^n \cdot \cos{((n-2k) \alpha)}$$

For odd $n$:

$$sin^n \alpha = \frac1{2^{n-1}} \cdot \sum_{k=0}^{\frac{n-1}{2}} (-1)^{\frac {n-1} 2 -k} \cdot C_k^n \cdot \sin{((n-2k) \alpha)}$$

$$cos^n \alpha = \frac1{2^{n-1}} \cdot \sum_{k=0}^{\frac{n-1}{2}} C_k^n \cdot \cos{((n-2k) \alpha)}$$


\subsubsection{Таблица интегралов}
\newcommand{\dx}{\hspace{2pt}dx}
\newcommand{\dd}[1]{\hspace{2pt}d#1}


\begin{minipage}{0.48\textwidth}
	\begin{multicols}{2}
		$$\int \frac{1}{x}\dx = \ln |x|$$
		
		$$\int a^x \dx = \frac{1}{\ln a} a^x$$

		$$\int \ln x \dx = x \ln x - x$$

		$$\int \tan x \dx = \ln |\sec x|$$

		$$\int \sec x \dx = \ln |\sec x + \tan x|$$
		
		$$\int \sec^2 x \dx = \tan x$$

		$$\int \sec x \tan x \dx = \sec x$$
		
		$$\int \frac{a}{a^2+x^2}\dx = \tan^{-1}\frac{x}{a}$$

		$$\int \frac{a}{a^2-x^2}\dx = \frac{1}{2}\ln\left|\frac{x+a}{x-a}\right|$$
		
		$$\int \frac{1}{\sqrt{a^2-x^2}} \dx = \sin^{-1} \frac{x}{a}$$

		$$\int \frac{a}{x \sqrt{x^2-a^2}} \dx = \sec^{-1} \frac{x}{a}$$
		
		$$\int \frac{1}{\sqrt{x^2-a^2}} \dx = \cosh^{-1} \frac{x}{a} = \nonumber \ln (x+\sqrt{x^2-a^2}) $$

		$$\int \frac{1}{\sqrt{x^2+a^2}} \dx = \sinh^{-1} \frac{x}{a} =\nonumber \ln (x+\sqrt{x^2+a^2})$$
		
		$$\int \!\sqrt {-{x}^{2}+1}\,{\rm d}x=\frac{1}{2}\,x\sqrt {-{x}^{2}+1}+\frac{1}{2}\,\arcsin \left( x \right) $$
	\end{multicols}
\end{minipage}

\subsubsection{Таблица производных}

\begin{minipage}{0.48\textwidth}
	\begin{multicols}{2}
		$$\arcsin(x)' = \frac{1}{\sqrt{1-x^2}}$$
		
		$$\arccos(x)' = \frac{-1}{\sqrt{1-x^2}}$$

		$$\arctan(x)' = \frac{1}{1x^2}$$
		
		$$\arccot(x)' = \frac{-1}{1+x^2}$$

		$$\log_a(x)' = \frac{1}{x \cdot \ln(a)}$$
		
		$$(a^x)' = a^x \cdot \ln(a)$$
	\end{multicols}
\end{minipage}

\subsubsection{Линейные уравнения по модулю}

Equation: $a \cdot x = b\ \pmod n$. We should divide $b$, $a$, $n$ on $\gcd {(a, n)}$.

$$a' \cdot x = b'\ \pmod n',\ x_i = x' + i \cdot n'\ \pmod n,\ i = 0..\gcd{(a, n)} - 1$$

\subsubsection{Числа Каталана}

Catalan number: 

$$C_n = \sum_{k = 0}^{n-1} C_k \cdot C_{n-1-k},\ C_n = \frac{1}{n+1} \cdot C_{2n}^{n}$$

$n$-th Catalan number is:
\begin {itemize}
	\item number of binary trees with $n+1$ unnumbered leaves;
	\item number of triangulations of convex $n+2$ polygon;
	\item number of connections $2n$ points on a circle with non-intersected chords;
	\item number of non-isomorphic full binary trees with $n$ vertices interior (that have at least one son);
	\item number of permutations of length $n$, which can be sorted by stack (not such indexes: $i < j < k$ and $a_k < a_i < a_j$);
	\item number of continuous partitions of set of $n$ elements;
	\item number of covering $n$-ladder with $n$ rectangles.
\end {itemize}
 
\subsubsection{Числа Стирлинга}

\begin {enumerate} [I]
	\item Number of permutations length $n$ with $k$ cycles.
	$$|s_{n, k}| = |s_{n-1, k-1}| + (n-1) \cdot |s_{n-1, k}|, s_{n, 1} = s_{n, n} = 1$$
	\item The number of ways to partition a set of n objects into k non-empty subsets.
	$$|S_{n, k}| = |S_{n-1, k-1}| + k \cdot |S_{n-1, k}|,\ S_{n, k} = \frac{1}{k!} \cdot \sum_{j=0}^{k} (-1)^{k-j}\cdot C_j^k \cdot j^n$$
\end {enumerate}

\subsubsection{Лемма Бернсайда}

Number of invariant objects: $$N = \frac{1}{\abs{|G|}}\sum_{\pi \in G} k^{C(\pi)}$$ where $G$ is set of all invariant permutations, 
$k$ is given number of colors, $C(\pi)$ is number of cycles in $\pi$


\subsubsection {Бином. коэффициенты}

\begin{minipage}{0.48\textwidth}
	\begin{multicols}{2}
		$$\sum_{k=0}^{n} C_{n}^{k} = 2^n$$
		
		$$\sum_{m=0}^{n} C_{m}^{k} = C_{n+1}^{k+1}$$

		$$\sum_{k=0}^{m} C_{n+k}^{k} = C_{n+m+1}^{m}$$
		
		$$\sum_{k=0}^{n} (C_{n}^{k})^2 = C_{2n}^{n}$$

		$$\sum_{k=1}^{n} (C_{n}^{k}) \cdot i = n \cdot 2 ^ {n-1}$$

		$$\sum_{k=0}^{n} (C_{n-k}^{k})^2 = F_{n+1}$$
	\end{multicols}
\end{minipage}

\subsubsection {Формулы}

\begin{minipage}{0.48\textwidth}
	\begin{multicols}{2}
	Geometric progression: $b_n = b_1 q^{n - 1}$
	$$
	S_n = \frac{b_1 (1 - q^n)}{1 - q},\textit { if } q \neq 1;\ S = \frac{b_1}{1 - q},\textit { if } \abs{q} < 1
	$$

	Pick's formula: 
	$$S_{polygon} = interior + \frac{boundary}{2} - 1$$

	Bell Numbers:
	$$B_{n+1} = \sum_{k=0}^{n}{C_k^n \cdot B_k} = \sum_{k=1}^{n}{S_{n,k}}$$
	
	Chinese theorem:
	$$x = \sum_{i=1}^{i=n}{q_i \cdot M_i \cdot M^{-1}_i}$$
	where $\ M_i = \frac{M}{p_i},\ M^{-1}_i$--- mult inverse of $M_i \bmod p_i$.

	Finding all the inverses in the ring modulo p for prime p:
	$$r[i] = - \floor {\frac{p}{i}} \cdot r[p\ \bmod\ i]\ \pmod p$$

	Sums of the series: 

	$$\sum_{k=0}^{n} k^2 = {\frac{n(n+1)(2n+1)}{6}}$$

	$$\sum_{k=0}^{n} k^3 = {\frac{n^2(n+1)^2}{4}}$$

	$$\sum_{k=0}^{n} k\cdot x^k = \frac{x - (n+1) \cdot x^{n+1} + n \cdot x^{n+2}}{(x-1)^2}$$

	Square Root: $$x_i = \frac{x_i + \frac{n}{x_i}}{2}$$

	Euler Formulas:
	\begin {itemize}
		\item If $G <V, E>$ --- planar graph with $F$ faces then $V - E + F = 2$.
		\item If $G<V, E>$ --- planar graph with $V \ge 3$ then $E \le 3V-6$.
	\end {itemize}
	\end{multicols}
	Simpson formula: 
	$$\int_a^b f(x) dx = \frac{b-a}{6} \cdot \left (f(a) + 4 \cdot f \left (\frac{a+b}{2} \right) + f(b) \right )$$

	Теорема Люка ($p$ --- простое, $n_i$, $m_i$ --- цифры в $p$-ой системе счисления):
	$$C_{n}^{m}\equiv \prod_{i=0}^{k-1}C_{n_i}^{m_i} \pmod p$$

	Критерий Эйлера ($p$ --- нечётное простое, $\gcd(p, a) = 1$):
	$$
	a^{\frac{p-1}{2}} \equiv
	\begin{cases}
	1 \pmod p, & \text{если } \exists x:\ x^2 \equiv a \pmod p \text{ (квадратичный вычет)},\\
	-1 \pmod p, & \text{иначе}.
	\end{cases}
	$$

\end{minipage}

\subsubsection{Пифагоровы тройки}

\begin{minipage}{0.48\textwidth}
    Примитивные тройки: $a=m^{2}-n^{2}, b=2mn, c=m^{2}+n^{2}$

    Все тройки: $a=k\cdot (m^{2}-n^{2}), b=k\cdot (2mn), c=k\cdot (m^{2}+n^{2})$, где $m, n, k$ — натуральные числа, $m > n$, $m - n$ нечётно, $gcd(m, n) = 1$
\end{minipage}